\section{Motivation}\label{sec:motivation}

In-context learning (ICL) recasts segmentation as few-shot matching:
given a target scan and a small set of (image, mask) context pairs
illustrating the structure of interest, a model predicts a mask for the
target with no task-specific fine-tuning. This removes the need to
retrain for every new anatomical class -- valuable for structures that
are rare, newly defined, or annotated at only a handful of sites.

Most ICL segmentation models were designed and validated in 2D. Moving to
native 3D volumes is attractive -- CT/MRI are inherently volumetric, and
slice-wise processing discards through-plane context needed to resolve
thin or elongated structures -- but it sharpens an already difficult
resolution/compute trade-off: dense cross-attention between a target
volume and its context volumes scales with the cube of side length, so a
design affordable at $64^3$ becomes impractical at $256^3$.

This trade-off is not one problem but three, and this thesis addresses
each with its own mechanism:

\begin{enumerate}[leftmargin=1.6em,itemsep=0.2em]
  \item \textbf{How should image and label evidence be combined?}
    Concatenating image and mask channels before a shared encoder, or
    fusing them once through a single operation, commits to one fusion
    point before most of the network's capacity is applied.
    \autoref{sec:methodology} proposes \emph{bi-axial image--label
    attention}, keeping the two streams separate and letting them
    exchange information repeatedly.
  \item \textbf{Where should compute be spent?} Re-scanning an entire
    volume with a sliding window at every processing step wastes compute
    on regions already known not to contain the structure of interest.
    \autoref{sec:methodology} proposes a \emph{coarse-to-fine cascade}
    that restricts fine processing to the region a coarse pass has
    already localized, connected across levels by a query prior and
    register tokens that carry a gradient path between resolutions.
  \item \textbf{Where does training signal come from beyond limited
    annotation?} Real (image, mask) pairs are limited to the classes and
    appearance distribution covered by the training cohort.
    \autoref{sec:methodology} proposes a \emph{synthetic task generator}
    that repaints real anatomical geometry with novel, controllable
    appearance, producing label-perfect training tasks beyond the
    annotated cohort.
\end{enumerate}

\paragraph{Contributions.}
\begin{itemize}[leftmargin=1.4em,itemsep=0.1em]
  \item A bi-axial image--label attention mechanism for in-context 3D
    segmentation, contrasted against early-concatenation
    (Medverse~\cite{huMedverseUniversalModel2025a}) and one-shot fusion
    (Iris~\cite{gaoShowSegmentUniversal2025}) baselines.
  \item A coarse-to-fine cascade in which the query prior,
    region-restricted fine processing, and cross-level register training
    are evaluated against the sliding-window-over-the-whole-volume design
    used by prior work.
  \item A GMM-based synthetic task generator with a texture-noise variant
    that better matches real tissue's spatial autocorrelation, evaluated
    against real-only and i.i.d.-noise training.
  \item Evaluation on TotalSegmentator CT and a range of
    out-of-distribution cohorts and modalities, reporting where each
    mechanism helps and where it does not.
\end{itemize}

\paragraph{Outline.} \autoref{sec:fundamentals} introduces in-context
segmentation, the volumetric encoder and attention machinery the method
builds on, and synthetic-data background, then positions UniverSeg, Iris,
and Medverse against the three axes above. \autoref{sec:methodology}
details the model. \autoref{sec:results} evaluates each axis in turn.
\autoref{sec:discussion} synthesizes what the evidence supports and what
remains open, and \autoref{sec:conclusion} summarizes.
